\section{Background and Related Work}
\label{sec:bg}

Let an ordered flow sequence be \(\mathbf{X}=\{x_t\}_{t=1}^{T}\), where each \(x_t\in\mathbb{R}^{d}\) is a normalized feature vector. Using window size \(W\) and stride \(r\), we construct window \(i\) as
\begin{equation}
\mathbf{W}_i=[x_{(i-1)r+1},\ldots,x_{(i-1)r+W}] \in \mathbb{R}^{W\times d}.
\end{equation}
A detector maps each window to a scalar anomaly score:
\begin{equation}
s_i=f(\mathbf{W}_i).
\end{equation}
An alarm is raised when \(s_i\ge\tau\), where \(\tau\) is calibrated from independent benign windows.

For evaluation, a window is labeled anomalous if its internal attack-record ratio is at least \(\rho\) (here \(\rho=0.15\)); otherwise it is benign. Given fixed \(\tau\), we report confusion-matrix counts \((\mathrm{TP},\mathrm{FP},\mathrm{TN},\mathrm{FN})\), precision, recall, and F1, along with AUC and AP for score ranking. Because deployment cost is dominated by benign alarms, we also report observed benign false-positive rate on test benign windows:
\begin{equation}
\mathrm{FPR}_{\mathrm{benign,test}}=\frac{\mathrm{FP}}{\mathrm{FP}+\mathrm{TN}}.
\end{equation}

The training and calibration philosophy is benign-centric. Scaler fitting, model training, and TAD reference construction use only model-train benign windows; threshold calibration uses only independent-calibration benign windows. Attack labels are not used to select thresholds, tune post-hoc operating points on test data, or train the unsupervised detector.

This formulation targets deployment-oriented intrusion detection rather than supervised attack-type classification. Attack labels are used only after scoring to compute evaluation metrics and attack-type breakdowns. As a result, the paper emphasizes calibrated alarm behavior and auditable evidence for raised alarms, not semantic multiclass attack recognition.

\subsection{Related Work Context}

\textbf{(Q1) How should low-noise IDS be evaluated for operations?}
From an applied-security viewpoint, threshold policy is a first-class design choice rather than a minor post-processing step. One-class and online IDS methods typically require explicit threshold setting and post-deployment benign-alert monitoring \cite{scholkopf2001ocsvm,mirsky2018kitsune,ruff2018deep}. Recent security analytics surveys also emphasize that time-series anomaly systems must be interpreted through deployment constraints, not ranking alone \cite{landauer2025timeseriessecurity}. This motivates our distinction between \emph{target FPR} (calibration rule) and \emph{observed test benign FPR} (operational outcome), which we report jointly with AUC/AP.

\textbf{(Q2) What are representative unsupervised NIDS baselines?}
Representative low-label baselines include classical detectors (IsolationForest, OneClassSVM) \cite{liu2008iforest,scholkopf2001ocsvm}, deep one-class/reconstruction methods (DeepSVDD, AnoGAN/f-AnoGAN) \cite{ruff2018deep,schlegl2017anogan,schlegl2019fanogan}, and adversarial anomaly detectors such as GANomaly \cite{akcay2018ganomaly}. Recent unsupervised network-centric GAN formulations further improve anomaly scoring quality on modern flow settings \cite{araujo2023unsupervised,ruffo2025ganipflow}. Our baseline set follows this representative scope under strict protocol alignment, rather than claiming exhaustive SOTA coverage.

\textbf{(Q3) What is the role of modern temporal/Transformer models?}
TCN backbones remain strong for efficient temporal modeling \cite{bai2018tcn}. In adversarial settings, WGAN and WGAN-GP provide critic-side training stability compared with vanilla GAN objectives \cite{goodfellow2014gan,arjovsky2017wgan,gulrajani2017wgangp}. Transformer-style and modern time-series models (TranAD, Anomaly Transformer, TimesNet, DLinear) have shown strong ranking performance in anomaly tasks \cite{tuli2022tranad,xu2022anomalytransformer,wu2023timesnet,zeng2023dlinear}, and recent model surveys further document the rapid evolution of deep time-series backbones \cite{wang2024tssurvey}. Recent flow-based/transformer NIDS studies continue this trend \cite{manocchio2024flowtransformer,long2024transformercloud,xi2024multiscale}. However, strong ranking does not by itself guarantee stable benign alarm rates at a fixed threshold, which is exactly the gap our strict calibration protocol targets.

\textbf{(Q4) What can XAI claim boundaries look like in IDS?}
Attribution tools such as Integrated Gradients and SHAP are increasingly used for IDS auditing and analyst support \cite{sundararajan2017axiomatic,lundberg2017shap,rahman2025explaining}. Domain-specific anomaly/XAI work in cyber-physical environments is also growing \cite{aslam2024improvedaeics}. At the same time, attention weights should not be treated as faithful explanation by default \cite{jain2019attention}. Our auditability claims therefore remain conservative: attribution-guided masking is presented as empirical decision-trace evidence, not proof of causal explanation correctness.

Dataset-wise, CICIDS2017 is the primary evidence chain in this paper \cite{sharafaldin2018cicids}. SWaT and UNSW-NB15 are used only for bounded supplementary/stress-test interpretation under the same strict philosophy \cite{goh2016swat,moustafa2015unsw}. We do not use these supplementary datasets to claim comprehensive cross-dataset generalization.
